\hmsubsection{Actuator Control}{FS}{}\label{sec:actuator-control}

The actuator control layer is the firmware running on the OpenRB-150.
Its job is to take goal positions from the host, drive the Dynamixel
servos, watch each motor, and report success or failure back to the
host. The firmware does not do any kinematics. The inverse kinematics
solver runs on the Pi (Section~\ref{sec:ik-solver}). The firmware
only sees encoder steps.

\subsubsection*{Firmware structure}

The firmware uses the Arduino \texttt{setup()} and \texttt{loop()}
pattern.
At \texttt{setup()}, the firmware opens the USB serial at 115\,200~bps,
opens the Dynamixel bus at 115\,200~bps using the Dynamixel2Arduino library,
and pings each of the five motors. If a motor does not respond, the
firmware sends \texttt{ERR INIT} to the host and goes into the
\texttt{FAULT} state. New commands are rejected until the fault is
cleared.

The state machine has three states: \texttt{INIT}, \texttt{READY},
\texttt{FAULT}. \texttt{READY} is the working state. The firmware
moves to \texttt{FAULT} when a motor reports a hardware error or a
motion times out. The host clears \texttt{FAULT} by sending a
\texttt{SCAN} command, which also moves the arm back to the scan pose.

The \texttt{loop()} function does three things on every pass. It
reads zero or more bytes from USB serial and appends them to a line
buffer. When the buffer has a newline, the line is parsed. If a
motion is in progress, it polls the motor positions and updates the
status. None of the three steps block, so the firmware can react to a
\texttt{STOP} from the host even while the arm is moving.

\subsubsection*{Motion execution}

A \texttt{MOVE} command from the host contains one goal position per
motor. The firmware first checks each value against the safe operating
range for that motor (Section~\ref{sec:joint-limits}). If any value
is out of range, the firmware sends \texttt{ERR OUT\_OF\_RANGE} and
the motion is rejected.

If all five values are valid, the firmware writes them to the motors
in one \texttt{syncWrite()} call. This sends one broadcast packet on
the Dynamixel bus instead of five per-motor packets, which cuts the
bus time from about 5~ms to about 1~ms. All five motors start moving
within roughly the same TTL slot.

After the write, the firmware sends \texttt{BUSY} to the host and
starts polling. It reads the present position of each motor at about
50~Hz and computes the difference from the goal. When every motor is
within the tolerance band (default 8~steps, about 0.7~degrees on the
XM430 family), the firmware sends \texttt{READY} and goes back to
reading commands.

\subsubsection*{Gripper handling}

The gripper motor (M5, XM430-W210-T) is the same model as the base
and elbow motors and lives on the same TTL bus. The firmware treats it
identically. The host does not send raw goal positions for the
gripper. It sends \texttt{GRIP open} or \texttt{GRIP close} instead.
The firmware turns each into the matching pre-calibrated goal
position. The open and closed positions are constants in the firmware,
set once during bring-up so the closed position grips a 50~mm ball
firmly without going over the torque limit. Using named tokens
instead of raw step values means the host does not need to know how
the gripper is tuned. The gripper can be re-tuned without changes to
the host.

\subsubsection*{Current-limit command}

The firmware supports a \texttt{LIMIT mN value} command that sets the
current limit of one motor through the Dynamixel \emph{Goal Current}
register. This mechanism was tested on the elbow motor (M3) with a
300~mA limit while the arm was parked at the scan pose, as described
in Section~\ref{sec:thermal}. Integration testing showed that the reduced limit did
not lower the thermal load sufficiently, so scan-pose current limiting
was removed from the deployed host behaviour. The firmware therefore
does not apply an automatic M3 current limit when the arm returns to
the scan pose; the deployed mitigation for idle scan-pose heating is
instead the mechanically supported scan posture described in
Section~\ref{sec:thermal}. This does not remove thermal loading during
movement. The same register-access mechanism remains available for
actuator diagnostics and is used for the gripper current cap in the
adaptive grip algorithm below.

\subsubsection*{Fault handling}

Several kinds of faults can happen during operation.
A motor can report a hardware error flag (overload, overheating,
communication). The firmware reads the error register of each motor
at about 5~Hz during motion. If any flag is set, the firmware sends
\texttt{ERR OVERLOAD} or \texttt{ERR OVERHEAT} to the host.

A motion can take longer than the cycle timeout. The firmware sends
\texttt{ERR TIMEOUT} and disables torque on all motors as a
precaution.

In both cases the firmware moves to \texttt{FAULT}. The host receives
the error token and aborts the cycle through the supervisor.
The firmware stays in
\texttt{FAULT} until the host sends \texttt{SCAN}, which clears the
fault and returns the arm to a known pose.


\hmsubsubsection{Adaptive Grip Algorithm}{UMA}{}

The firmware's simple \texttt{GRIP close} command drives the gripper
to a fixed closed position, but it cannot distinguish between
successfully gripping a ball and closing on air. Without feedback,
the system would attempt to transport a ball that was never picked
up, wasting a cycle and potentially desynchronising the state
machine. To address this, the host implements an adaptive grip
algorithm that exploits the Dynamixel protocol's real-time access
to the motor's present position, present load, and present current
registers. This approach constitutes a form of impedance-based
force sensing that requires no external sensor hardware and is
specific to smart actuators such as the Dynamixel series.

\textbf{Problem.} A conventional open/close gripper command
provides no information about whether an object is actually held.
If the ball has rolled away, been knocked aside during approach,
or was never within reach, the claw closes fully without
resistance. Without grip verification, the system proceeds to the
transport phase with an empty gripper, which wastes time and may
trigger spurious sort failures.

\textbf{Algorithm.} The adaptive grip proceeds in five stages:

\begin{enumerate}
    \item \textbf{Current cap.} The M5 current limit is set to
    50~mA via the Dynamixel \emph{Goal Current} register. This
    value is low enough to detect resistance from a 50~mm ball
    but gentle enough to protect the 3D-printed claw mechanism
    from damage.

    \item \textbf{Slow motion profile.} The profile velocity is
    set to 80 and the profile acceleration to 20, producing a
    controlled closing motion that gives the polling loop
    sufficient time to detect contact.

    \item \textbf{Open.} The claw is commanded to the open
    position (2745~steps).

    \item \textbf{Incremental close with sensor polling.} The
    claw advances toward the closed position (3350~steps) in
    increments of 50~steps. At each increment, the host polls
    the motor at 50~ms intervals for three sensor channels:
    \begin{itemize}
        \item \emph{Load}: contact is detected if the present
              load exceeds 5~units.
        \item \emph{Current}: contact is detected if the
              present current exceeds 20~mA (40\% of the
              50~mA current cap).
        \item \emph{Stall}: contact is detected if the motor
              position advances by fewer than 8~steps for two
              consecutive readings.
    \end{itemize}
    The loop terminates on the first contact detection or after
    a 15-second timeout.

    \item \textbf{Secure close.} Once contact is detected, the
    claw is commanded to the full closed position (3350~steps)
    to squeeze around the ball, and the system waits for the
    claw to settle.
\end{enumerate}

\textbf{Verification.} After the secure close, the final motor
position, load, and current are read and evaluated against a
two-tier threshold:

\begin{itemize}
    \item \textbf{Strongly blocked.} The claw stopped more than
    30~steps from the closed position. Position alone is
    sufficient to confirm grip, provided some resistance is
    detected (load~$> 0$ or current~$\geq 5$~mA).

    \item \textbf{Minimally blocked.} The claw stopped at least
    5~steps from the closed position. Grip is confirmed only if
    the load~$\geq 15$ or the current~$\geq 20$~mA.

    \item \textbf{Zero resistance.} If both load and current are
    at or near zero, the grip is rejected immediately regardless
    of position --- the claw closed on air.

    \item \textbf{At end-stop.} If the claw is within 5~steps of
    the closed position, the grip is always rejected --- no ball
    was present.
\end{itemize}

If the adaptive grip returns an inconclusive result, a secondary
verification function reads the sensor feedback once more and
applies the same two-tier check. If both attempts fail, the claw
opens and the ball is abandoned; the system returns to the scan
position for the next detection cycle.

\textbf{Why this approach works.} By capping the current limit to
50~mA, the system amplifies the difference between closing on air
(near-zero resistance) and closing on a ball (measurable resistance
that stalls the motor before it reaches the closed position). The
Dynamixel XM430-W210-T exposes its present position, present load,
and present current as readable registers via Protocol~2.0, making
this form of impedance-based sensing possible without any additional
hardware. Conventional hobby servos do not provide current or load
feedback, so this approach is specific to smart actuator platforms.

The adaptive grip constants are defined in
\texttt{src/config/arm.py} and can be adjusted without modifying
the control logic. The algorithm was validated using the
calibration script \texttt{02b\_claw\_grip\_test.py}, which mirrors
the production algorithm and provides per-attempt sensor logging
for threshold tuning.

\begin{table}[H]
    \centering
    \renewcommand{\arraystretch}{1.3}
    \begin{tabular}{|l|r|p{5.5cm}|}
        \hline
        \textbf{Constant} & \textbf{Value} &
        \textbf{Description} \\
        \hline
        \texttt{CLAW\_OPEN\_POS}   & 2745 &
        Gripper fully open (steps) \\
        \hline
        \texttt{CLAW\_CLOSED\_POS} & 3350 &
        Gripper closed limit (steps) \\
        \hline
        \texttt{GRIP\_CURRENT\_LIMIT} & 50~mA &
        M5 current cap during grip \\
        \hline
        \texttt{GRIP\_PROFILE\_VEL} & 80 &
        Closing velocity \\
        \hline
        \texttt{GRIP\_PROFILE\_ACC} & 20 &
        Closing acceleration \\
        \hline
        \texttt{GRIP\_POLL\_INTERVAL} & 0.05~s &
        Sensor polling rate \\
        \hline
        \texttt{GRIP\_TIMEOUT} & 15.0~s &
        Maximum grip duration \\
        \hline
        \texttt{GRIP\_LOAD\_DETECT} & 5 &
        Light contact threshold \\
        \hline
        \texttt{GRIP\_LOAD\_THRESHOLD} & 15 &
        Firm grip confirmation \\
        \hline
        \texttt{GRIP\_POSITION\_STALL} & 8 &
        Stall detection (steps) \\
        \hline
        \texttt{GRIP\_EXTRA\_CLOSE} & 50 &
        Step increment per iteration \\
        \hline
    \end{tabular}
    \caption{Adaptive grip constants defined in
    \texttt{src/config/arm.py}. All values were determined
    empirically during claw calibration.}
\end{table}
